\subsection{Behavior of the product with respect to the usual operations}

PROPOSITION 5. --- Let \((\lambda_{\alpha})_{\alpha\in\mathbb{A}}\) be a family of positive measures on T, directed for the relation \(\leqslant\) , admitting in \(\mathcal{M}(\mathrm{T})\) a supremum \(\lambda\) . For a positive numerical function g to be locally \(\lambda\) -integrable, it is necessary and sufficient that g be locally \(\lambda_{\alpha}\) -integrable for every \(\alpha\in A\) and that the family \((g\cdot\lambda_{\alpha})\) be bounded above in \(\mathcal{M}(\mathrm{T})\) ; in this case,

\[
g \cdot \lambda = \sup _ {\alpha \in \mathrm{A}} g \cdot \lambda_ {\alpha}.
\]

It is clear that the condition is necessary. Conversely, suppose that \(g\) is locally integrable for each of the measures \(\lambda_{\alpha}\) and that the family \((g \cdot \lambda_{\alpha})_{\alpha \in \mathbb{A}}\) is bounded above; denote its supremum by \(\lambda'\) . The function \(g\) is then \(\lambda\) -measurable (§1, No.~4, Cor. 2 of Prop. 11); moreover, for every function \(h \in \mathcal{K}_{+}(\mathbf{T})\) ,

\[
\int^ {\bullet} (h g) d \lambda = \sup _ {\alpha \in A} \int^ {\bullet} (h g) d \lambda_ {\alpha} = \sup _ {\alpha \in A} \int^ {\bullet} h d (g \cdot \lambda_ {\alpha}) = \int^ {\bullet} h d \lambda^ {\prime}
\]

(§1, No.~4, Prop. 11). This implies first of all that the first member is finite for any \(h\) , so that \(g\) is locally \(\lambda\) -integrable; the symbol \(\int^{\bullet}\) may therefore be replaced by \(\int\) , and the formula may be written \(\int h d(g \cdot \lambda) = \int h d\lambda'\) . It follows that \(g \cdot \lambda = \lambda'\) , and this completes the proof.

COROLLARY. --- Suppose that \(\mu\) is the sum of a family \((\mu_{\alpha})_{\alpha\in\mathrm{A}}\) of measures on T. For a positive numerical function g defined on T to be locally \(\mu\) -integrable, it is necessary and sufficient that g be locally \(\mu_{\alpha}\) -integrable for every \(\alpha\in A\) and that the family \((g\cdot\mu_{\alpha})_{\alpha\in\mathrm{A}}\) be summable. In this case,

\[
g \cdot \mu = \sum_ {\alpha \in A} g \cdot \mu_ {\alpha}.\tag{7}
\]

Let \((g_{\alpha})_{\alpha \in A}\) be a family of \(\mu\) -measurable positive functions defined on T. Let \(S_{\alpha}\) be the set of \(t \in T\) such that \(g_{\alpha}(t) \neq 0\) . We shall say that the family \((g_{\alpha})\) is locally countable if the family \((S_{\alpha})\) is locally countable (Ch. IV, §5, No.~9); this amounts to saying that, for every compact set K in T, the set of \(\alpha \in A\) such that \(g_{\alpha}|K\) is not zero is countable.

PROPOSITION 6. --- Let \((g_{\alpha})_{\alpha \in A}\) be a locally countable family of locally \(\mu\) -integrable positive numerical functions defined on T. In order that the function \(g = \sum_{\alpha \in A} g_{\alpha}\) be locally \(\mu\) -integrable, it is necessary and sufficient that the family of measures \((g_{\alpha} \cdot \mu)_{\alpha \in A}\) be summable, in which case

\[
g \cdot \mu = \sum_ {\alpha \in A} g _ {\alpha} \cdot \mu .\tag{8}
\]

It is clear that \(g\) is \(\mu\) -measurable (Ch. IV, §5, No.~2, Prop. 4 and No.~4, Cor. 1 of Th. 2). For \(g\) to be locally \(\mu\) -integrable, it is therefore necessary and sufficient that \(\mu^{\bullet}(gf)\) be finite for every \(f \in \mathcal{K}_{+}(\mathrm{T})\) . Now, since the set of \(\alpha \in \mathrm{A}\) such that \(g_{\alpha}f \neq 0\) is countable, we have \(\mu^{\bullet}(gf) = \sum_{\alpha \in \mathrm{A}} \mu^{\bullet}(g_{\alpha}f)\) (§1, No.~1, Cor. of Prop. 2). Set \(\nu_{\alpha} = g_{\alpha} \cdot \mu\) ; the condition \(\mu^{\bullet}(gf) < +\infty\) is equivalent to the condition \(\sum_{\alpha \in \mathrm{A}} \nu_{\alpha}(f) < +\infty\) : in other words, \(g\) is locally \(\mu\) -integrable if and only if the family \((\nu_{\alpha})\) is summable. Denoting the sum of this family by \(\nu\) , the preceding calculation yields the equality \(\nu(f) = \mu^{\bullet}(gf)\) , which is equivalent to (8).

COROLLARY. --- Let \((g_{\alpha})\) be a sequence of locally \(\mu\) -integrable numerical functions, such that the sequence of measures \(g_{n} \cdot \mu\) is increasing. In order that this sequence have an upper bound in the ordered vector space \(\mathcal{M}(T)\) of real measures on T, it is necessary and sufficient that the function \(g = \sup g_{n}\) be locally \(\mu\) -integrable; the supremum in \(\mathcal{M}(T)\) of the sequence \((g_{n} \cdot \mu)\) is then the measure \(g \cdot \mu\) .

It suffices to apply Prop. 6 to the functions (positive locally almost everywhere) \(g_n' = g_{n+1} - g_n\) .

PROPOSITION 7. --- Let X be a locally compact space that is countable at infinity, and let \(t \mapsto \lambda_t\) be a \(\mu\) -adequate mapping of T into \(\mathcal{M}_+(X)\) .

Let \(g\) be a positive numerical function defined on \(X\) , locally integrable for the measure \(\nu = \int \lambda_t d\mu(t)\) . The set of \(t \in T\) such that \(g\) is not locally \(\lambda_t\) -integrable is then locally negligible for \(\mu\) , the mapping \(t \mapsto g \cdot \lambda_t\) (defined locally \(\mu\) -almost everywhere) is \(\mu\) -adequate, and

\[
g \cdot \nu = \int (g \cdot \lambda_ {t}) d \mu (t).\tag{9}
\]

Let \((\mathrm{K}_{n})_{n\in\mathbf{N}}\) be an increasing sequence of compact subsets of X whose interiors cover X; if \(\eta\) is any positive measure on X, to say that g is locally \(\eta\) -integrable is equivalent to saying that \(g\varphi_{\mathrm{K}_{n}}\) is \(\eta\) -integrable for every n.~Now let \(\mathrm{H}_{n}\) be the set of \(t\in\mathrm{T}\) such that \(g\varphi_{\mathrm{K}_{n}}\) is not \(\lambda_{t}\) -integrable, and let \(\mathrm{H}=\bigcup_{n}\mathrm{H}_{n}\) ; since \(\mathrm{H}_{n}\) is locally \(\mu\) -negligible for all n (§3, No.~3, Th. 1), the same is true of H, which establishes the first assertion of the statement. Replacing \(\lambda_{t}\) by 0 for t in H (which does not change the measure \(\nu\) ), we can suppose that g is locally \(\lambda_{t}\) -integrable for every \(t\in\mathrm{T}\) . For every \(\nu\) -measurable positive function h defined on X, we have, by Prop. 3 and by Prop. 5 of §3, No.~2,

\[
\int^ {\bullet} h d (g \cdot \nu) = \int^ {\bullet} (g h) d \nu = \int^ {\bullet} d \mu (t) \int^ {\bullet} (g h) d \lambda_ {t} = \int^ {\bullet} d \mu (t) \int^ {\bullet} h d (g \cdot \lambda_ {t}).
\]

This formula and Prop. 5 of §3, No.~2 first show (on taking \(h \in \mathcal{K}_{+}(\mathbf{X})\) ) that the mapping \(t \mapsto g \cdot \lambda_{t}\) is scalarly essentially \(\mu\) -integrable, and that its integral is \(g \cdot \nu\) ; in other words, the relation (9) holds. Next, let us replace \(\mu\) by a positive measure \(\mu' \leqslant \mu\) , and let us take for \(h\) a positive lower semi-continuous function: it follows at once from these relations that \(t \mapsto g \cdot \lambda_{t}\) is \(\mu\) -adequate (§3, No.~1, Def. 1).

PROPOSITION 8. --- Let \(\theta\) be a complex measure on T, \(g_{1}\) a locally \(\theta\) -integrable complex function, \(\theta_{1}\) the measure \(g_{1} \cdot \theta\) . For a complex function \(g_{2}\) defined on T to be locally \(\theta_{1}\) -integrable, it is necessary and sufficient that \(g_{2} g_{1}\) be locally \(\theta\) -integrable, in which case

\[
g _ {2} \cdot \theta_ {1} = g _ {2} \cdot (g _ {1} \cdot \theta) = (g _ {2} g _ {1}) \cdot \theta\tag{10}
\]

(`associativity formula').

By the corollary of Prop. 4, to say that \(g_{2}\) is \(\theta_{1}\) -measurable is equivalent to saying that \(g_{2}g_{1}\) is \(\theta\) -measurable. Let us suppose this condition to be satisfied. For every function \(f \in \mathcal{K}_{+}(T)\) we have, by Propositions 2 and 3,

\[
\int^ {\bullet} | g _ {2} | f d | \theta_ {1} | = \int^ {\bullet} | g _ {2} | f | g _ {1} | d | \theta | = \int^ {\bullet} | g _ {2} g _ {1} | f d | \theta |.
\]

To say that \(g_{2}\) is locally \(\theta_{1}\) -integrable is therefore equivalent to saying that \(g_{2}g_{1}\) is locally \(\theta\) -integrable. Assuming this condition to be satisfied, by Th. 1 we have

\[
\int f d \left(g _ {2} \cdot \theta_ {1}\right) = \int f g _ {2} d \theta_ {1} = \int f g _ {2} g _ {1} d \theta = \int f d \left(g _ {2} g _ {1} \cdot \theta\right),
\]

a formula equivalent to (10).
% semantic-fragments: legacy-input-seam
\relax{}
